% Original fixed-tier evaluation fixture for Set P (not a real paper).
% Tests comment traps, multi-label cleveref, inner equation in proof, and diverse kinds.
\documentclass{article}
\usepackage{amsmath,amssymb,amsthm,cleveref}

\newtheorem{theorem}{Theorem}[section]
\newtheorem{lemma}[theorem]{Lemma}
\newtheorem{corollary}[theorem]{Corollary}
\newtheorem{conjecture}[theorem]{Conjecture}
\theoremstyle{definition}
\newtheorem{assumption}[theorem]{Assumption}
\newtheorem{example}[theorem]{Example}
\theoremstyle{remark}
\newtheorem{remark}[theorem]{Remark}

\begin{document}

\section{Subadditive cocycles}
\label{sec:cocycles}

\begin{assumption}
\label{assump:integrable}
Let $(\Omega, \mathcal{F}, \mathbb{P})$ be a probability space equipped with an
ergodic measure-preserving map $T : \Omega \to \Omega$, and let
$M : \Omega \to \mathrm{GL}_d(\mathbb{R})$ be a measurable matrix-valued map
with positive entries satisfying $\log^+ \|M\| \in L^1(\mathbb{P})$.
\end{assumption}

\begin{lemma}
\label{lem:subadd}
Define the cocycle products $A_n(\omega) = M(T^{n-1}\omega) \cdots M(\omega)$
and $F_n(\omega) = \log \|A_n(\omega)\|$ for $n \ge 1$. Then for all integers
$m, n \ge 1$ and every $\omega \in \Omega$,
\begin{equation}
\label{eq:kingman-step}
F_{m+n}(\omega) \le F_m(\omega) + F_n(T^m \omega).
\end{equation}
\end{lemma}
\begin{proof}
Because $A_{m+n}(\omega) = A_n(T^m \omega) A_m(\omega)$, submultiplicativity
of the operator norm yields
$\|A_{m+n}(\omega)\| \le \|A_n(T^m \omega)\| \|A_m(\omega)\|$. Taking
logarithms gives the subadditive inequality.
\end{proof}

% \begin{theorem}
% \label{thm:commented-trap}
% Commented-out statement referencing Lemma~\ref{lem:subadd}.
% \end{theorem}

\section{Lyapunov exponent and examples}
\label{sec:lyapunov}

\begin{theorem}
\label{thm:lyapunov}
Under \cref{assump:integrable,lem:subadd}, there exists a constant
$\Lambda \in [-\infty, \infty)$ such that
\[
\lim_{n \to \infty} \frac{1}{n} F_n(\omega) = \Lambda
\qquad \text{for } \mathbb{P}\text{-almost every } \omega \in \Omega,
\]
and moreover $\Lambda = \inf_{n \ge 1} \frac{1}{n} \mathbb{E}[F_n]$.
\end{theorem}
\begin{proof}
Taking expectations in \eqref{eq:kingman-step} and using $T$-invariance of
$\mathbb{P}$ gives $\mathbb{E}[F_{m+n}] \le \mathbb{E}[F_m] + \mathbb{E}[F_n]$,
so $a_n = \mathbb{E}[F_n]$ is subadditive and $a_n / n \to \Lambda = \inf_{n \ge 1} a_n / n$.
% By Assumption~\ref{assump:integrable} we could stop early.
Iterating \eqref{eq:kingman-step} along blocks of length $k$ and applying
Birkhoff's ergodic theorem yields $\limsup_{n \to \infty} n^{-1} F_n(\omega) \le a_k / k$
almost surely; sending $k \to \infty$ establishes the upper bound, while the
standard maximal inequality supplies the matching almost-sure lower bound.
\end{proof}

\begin{corollary}
\label{cor:rate}
If in addition $\log^+ \|M^{-1}\| \in L^1(\mathbb{P})$, then the top
Lyapunov exponent $\Lambda$ is finite and obeys
$\Lambda \ge -\mathbb{E}[\log^+ \|M^{-1}\|]$.
% Suppose Assumption~\ref{assump:integrable} fails.
\end{corollary}
\begin{proof}
Since $1 = \|A_n(\omega)^{-1} A_n(\omega)\| \le \|A_n(\omega)^{-1}\| \|A_n(\omega)\|$,
we have $F_n(\omega) \ge -\sum_{j=0}^{n-1} \log^+ \|M(T^j \omega)^{-1}\|$.
Dividing by $n$, taking expectations, and applying Theorem~\ref{thm:lyapunov}
gives $\Lambda \ge -\mathbb{E}[\log^+ \|M^{-1}\|] > -\infty$.
\end{proof}

\begin{example}
\label{ex:diagonal}
If $M(\omega) = \operatorname{diag}(e^{X_1(\omega)}, \dots, e^{X_d(\omega)})$
for integrable random variables $X_1, \dots, X_d$, then the cocycle is
commutative and $\Lambda = \max_{1 \le i \le d} \mathbb{E}[X_i]$.
\end{example}

\begin{remark}
\label{rem:nonnegative}
When every entry of $M(\omega)$ is at least $1$, each product satisfies
$\|A_n(\omega)\| \ge 1$ and hence $\Lambda \ge 0$.
\end{remark}

\begin{conjecture}
\label{conj:analyticity}
For any positive measurable perturbation $N : \Omega \to \mathbb{R}^{d \times d}$
with $\log^+ \|N\| \in L^1(\mathbb{P})$, the top Lyapunov exponent
$\varepsilon \mapsto \Lambda(M + \varepsilon N)$ is real-analytic on an open
interval containing $[0, \infty)$.
\end{conjecture}

\end{document}
